\section*{Capítulo \ref{ch:g8:the-atmosphere} --- La atmósfera}

\begin{solution}{exo:g8:the-atmosphere:1}
(a) \omterm{def:g8:the-atmosphere:layers}{estratosfera}; (b) \omterm{def:g8:the-atmosphere:layers}{troposfera}; (c) \omterm{def:g8:the-atmosphere:layers}{mesosfera}; (d) \omterm{def:g8:the-atmosphere:layers}{termosfera}.
\end{solution}

\begin{solution}{exo:g8:the-atmosphere:2}
(a) \qty{15}{\degreeCelsius}. (b) Unos \qty{-56}{\degreeCelsius}.
(c) Unos \qty{-3}{\degreeCelsius}. (d) Unos \qty{-86}{\degreeCelsius}.
(e) De \qty{11}{km} a unos \qty{20}{km} (y a lo largo de unos pocos
kilómetros cerca de \qty{50}{km} y justo por encima de \qty{86}{km}).
\end{solution}

\begin{solution}{exo:g8:the-atmosphere:3}
(a) Falso: disminuye en la \omterm{def:g8:the-atmosphere:layers}{troposfera}, aumenta en la \omterm{def:g8:the-atmosphere:layers}{estratosfera},
disminuye en la \omterm{def:g8:the-atmosphere:layers}{mesosfera} y vuelve a aumentar en la \omterm{def:g8:the-atmosphere:layers}{termosfera}.
(b) Verdadero. (c) Falso: está en la \omterm{def:g8:the-atmosphere:layers}{estratosfera}. (d) Verdadero.
\end{solution}

\begin{solution}{exo:g8:the-atmosphere:4}
(a) La \omterm{def:g8:the-atmosphere:layers}{troposfera}. (b) La \omterm{def:g8:the-atmosphere:layers}{estratosfera} (en la atmósfera estándar, la
\omterm{def:g8:the-atmosphere:layers}{tropopausa} está a \qty{11}{km}). (c) La \omterm{def:g8:the-atmosphere:layers}{estratosfera}. (d) La \omterm{def:g8:the-atmosphere:layers}{mesosfera}.
\end{solution}

\begin{solution}{exo:g8:the-atmosphere:5}
(a) $20 - 6.5 \times 2 = 20 - 13 = \qty{7}{\degreeCelsius}$. (b) $20 -
6.5 \times 3 = 20 - 19.5 = \qty{0.5}{\degreeCelsius}$. (c) El aire tiene
que enfriarse 20 grados: $20 \div 6.5 \approx 3.1$, así que a unos
\qty{3.1}{km}.
\end{solution}

\begin{solution}{exo:g8:the-atmosphere:6}
(a) $3 - 0.5 = \qty{2.5}{km}$. (b) $6.5 \times 2.5 = 16.25$: unos 16
grados. (c) $18 - 16.25 \approx \qty{2}{\degreeCelsius}$.
\end{solution}

\begin{solution}{exo:g8:the-atmosphere:7}
(a) \qty{500}{hPa}, \qty{250}{hPa} y \qty{125}{hPa}. (b) Cerca del suelo
la regla funciona bien (\qty{500}{hPa} frente a \qty{506}{hPa}). Más
arriba da demasiado: \qty{250}{hPa} frente a \qty{227}{hPa}, y
\qty{125}{hPa} frente a \qty{96}{hPa}. La regla de la mitad solo es una
guía aproximada.
\end{solution}

\begin{solution}{exo:g8:the-atmosphere:8}
(a) $10~\text{hPa} = 10^1$~hPa; $10^3 \div 10^1 = 10^2$: cien veces menor.
(b) $0.01~\text{hPa} = 10^{-2}$~hPa. (c) $10^3 \div 10^{-2} =
10^5$: cien mil veces menor.
\end{solution}

\begin{solution}{exo:g8:the-atmosphere:9}
(a) $6 \div 1.5 = 4$ veces más ancho. (b) A medida que sube, el aire que lo
rodea presiona cada vez menos, así que el gas de dentro se expande y el
globo crece hasta que su envoltura se rompe. (c) Después de que el globo
estalle, el paracaídas frena la caída de la radiosonda, para que no hiera
a nadie ni dañe nada.
\end{solution}

\begin{solution}{exo:g8:the-atmosphere:10}
(a) $20 \div 7 \approx 2.9$ veces más gruesa. (b) Más cerca de la de los
\omterm{def:g3:continents-and-oceans:poles}{polos}: \qty{11}{km} está \qty{4}{km} por encima de \qty{7}{km}, pero
\qty{9}{km} por debajo de \qty{20}{km} (describe las \omterm{def:g6:latitude-longitude:latitude}{latitudes} medias, no
el \omterm{def:g3:continents-and-oceans:equator}{ecuador}). (c) Al principio, la \omterm{def:g8:the-atmosphere:layers}{tropopausa} está a unos \qty{7}{km}, por
debajo del avión: el avión está en la \omterm{def:g8:the-atmosphere:layers}{estratosfera}. Al final, la
\omterm{def:g8:the-atmosphere:layers}{tropopausa} está a unos \qty{20}{km}: el avión está en la \omterm{def:g8:the-atmosphere:layers}{troposfera}.
\end{solution}

\begin{solution}{exo:g8:the-atmosphere:11}
(a) $1 - 506 \div 1013 \approx 0.50$: \textbf{50~\%}. (b) $1 - 227 \div
1013 \approx 0.78$: \textbf{78~\%}. (c) $1 - 12 \div 1013 \approx 0.99$:
\textbf{99~\%}. (d) Alrededor del 1~\%.
\end{solution}

\begin{solution}{exo:g8:the-atmosphere:12}
(a) $6371 \div 15 \approx 425$~km por cada centímetro. (b) $100 \div 425
\approx 0.24$~cm, unos \qty{2.4}{mm}. (c) $11 \div 425 \approx
0.026$~cm, unos \qty{0.26}{mm}. (d) $0.26 \div 0.1 \approx 2.6$: dos o
tres hojas de papel.
\end{solution}

\begin{solution}{pb:g8:the-atmosphere:1}
\textbf{1.} $300 \times 10 = 3000$~m.

\textbf{2.} $6000 \div 300 = 20$ minutos; $\num{30000} \div 300 = 100$
minutos.

\textbf{3.} 100 minutos, es decir, 1~hora y 40~minutos.

\textbf{4.} A \qty{8}{km}: la \omterm{def:g8:the-atmosphere:layers}{troposfera} (la temperatura sigue bajando).
A \qty{30}{km}: la \omterm{def:g8:the-atmosphere:layers}{estratosfera}.

\textbf{5.} $15 - (-50) = 65$ grados.

\textbf{6.} $65 \div 10 = \qty{6.5}{\degreeCelsius}$ por kilómetro.

\textbf{7.} De \qty{12}{km} a \qty{20}{km} se mantiene en
\qty{-56}{\degreeCelsius}. El enfriamiento se detiene entre \qty{10}{km} y
\qty{12}{km}: la \omterm{def:g8:the-atmosphere:layers}{tropopausa} está a unos \qty{11}{km}.

\textbf{8.} Aumenta: \qty{-55}{\degreeCelsius} a \qty{22}{km},
\qty{-47}{\degreeCelsius} a \qty{30}{km}. El globo está en la
\omterm{def:g8:the-atmosphere:layers}{estratosfera}, calentada por el ozono, que absorbe la luz ultravioleta del
Sol.

\textbf{9.} $1013 \div 2 \approx \qty{507}{hPa}$; se alcanza entre
\qty{4}{km} (\qty{617}{hPa}) y \qty{6}{km} (\qty{472}{hPa}).

\textbf{10.} $1013 \div 4 \approx \qty{253}{hPa}$.

\textbf{11.} Entre \qty{10}{km} (\qty{265}{hPa}) y \qty{12}{km}
(\qty{194}{hPa}).

\textbf{12.} Desde \qty{10}{km}, la presión tiene que bajar $265 - 253 =
12$~hPa más, de los 71~hPa que pierde en 2~km: $12 \div 71 \times 2
\approx 0.34$~km. La presión es la cuarta parte de su valor al \omterm{def:g4:world-physical-map:altitude}{nivel del
mar} a unos \textbf{\qty{10.3}{km}}.
\end{solution}
